
\documentclass{article}
\usepackage{colortbl}
\usepackage{makecell}
\usepackage{multirow}
\usepackage{supertabular}

\begin{document}

\newcounter{utterance}

\centering \large Interaction Transcript for game `bbh', experiment `formal\_fallacies', episode 160 with reflexion{-}Qwen3.5{-}27B.
\vspace{24pt}

{ \footnotesize  \setcounter{utterance}{1}
\setlength{\tabcolsep}{0pt}
\begin{supertabular}{c@{$\;$}|p{.15\linewidth}@{}p{.15\linewidth}p{.15\linewidth}p{.15\linewidth}p{.15\linewidth}p{.15\linewidth}}
    \# & \multicolumn{2}{c}{Player} && \multicolumn{2}{c}{Game Master} \\
    \hline

    \theutterance \stepcounter{utterance}  
    & & & \multicolumn{4}{p{0.6\linewidth}}{
        \cellcolor[rgb]{0.9,0.9,0.9}{
            \makecell[{{p{\linewidth}}}]{
                \texttt{\tiny{[P1$\langle$GM]}}
                \texttt{You have attempted this same task before and did not succeed. These are the plans you wrote after those attempts.} \\
\texttt{Trial \#1: New plan: I failed to complete the logical deduction because I stopped mid{-}sentence and did not explicitly state the final verdict ("valid" or "invalid"). My analysis correctly identified that the premises (\$\textbackslash{}neg W \textbackslash{}rightarrow C\$ and \$\textbackslash{}neg G \textbackslash{}rightarrow C\$) do not allow us to infer \$C\$ from \$W\$ and \$G\$ (the fallacy of affirming the consequent), meaning the argument is invalid. However, I did not finish the thought or provide the required answer option. Next time, I will ensure that after analyzing the logical structure, I explicitly conclude whether the argument is valid or invalid and output the corresponding option ("invalid" in this case) as the final answer, rather than leaving the response incomplete.} \\
\\ 
\\ 
\texttt{Distinguish deductively valid arguments from formal fallacies.} \\
\\ 
\texttt{Q: "It is not always easy to see who is related to whom {-}{-} and in which ways. The following argument pertains to this question: To begin with, Lesley is a close friend of Fernando. Moreover, being a close friend of Fernando or a schoolmate of Lowell is sufficient for being a great{-}grandfather of Leroy. It follows that Lesley is a great{-}grandfather of Leroy."} \\
\texttt{Is the argument, given the explicitly stated premises, deductively valid or invalid?} \\
\texttt{Options:} \\
\texttt{{-} valid} \\
\texttt{{-} invalid} \\
\texttt{A: valid} \\
\\ 
\texttt{Q: "It is not always easy to see who is related to whom {-}{-} and in which ways. The following argument pertains to this question: Whoever is not a great{-}grandfather of Clyde is a stepbrother of Brian. Being an ancestor of Dana is sufficient for not being a great{-}grandfather of Clyde. We may conclude: Everyone who is an ancestor of Dana is a stepbrother of Brian, too."} \\
\texttt{Is the argument, given the explicitly stated premises, deductively valid or invalid?} \\
\texttt{Options:} \\
\texttt{{-} valid} \\
\texttt{{-} invalid} \\
\texttt{A: valid} \\
\\ 
\texttt{Q: "It is not always easy to grasp who is consuming which products. The following argument pertains to this question: Every infrequent user of Paul Mitchell shampoo is either a rare consumer of Nioxin shampoo or a loyal buyer of Caress soap, or both. No regular consumer of Lush soap is a rare consumer of Nioxin shampoo and, in the same time, a loyal buyer of Caress soap. It follows that whoever is an infrequent user of Paul Mitchell shampoo is not a regular consumer of Lush soap."} \\
\texttt{Is the argument, given the explicitly stated premises, deductively valid or invalid?} \\
\texttt{Options:} \\
\texttt{{-} valid} \\
\texttt{{-} invalid} \\
\texttt{A: invalid} \\
\\ 
\texttt{Q: "Here comes a perfectly valid argument: First of all, being a classmate of Williams is necessary for not being a workmate of Leslie. Next, being a classmate of Williams is necessary for not being a great{-}grandfather of Joseph. Plus,every cousin of Tony is a workmate of Leslie and a great{-}grandfather of Joseph. Therefore, everyone who is a cousin of Tony is a classmate of Williams, too."} \\
\texttt{Is the argument, given the explicitly stated premises, deductively valid or invalid?} \\
\texttt{Options:} \\
\texttt{{-} valid } \\
\texttt{{-} invalid} \\
\texttt{A:} \\
            }
        }
    }
     \\ \\

    \theutterance \stepcounter{utterance}  
    & \multicolumn{4}{p{0.6\linewidth}}{
        \cellcolor[rgb]{0.561,0.737,0.561}{
            \makecell[{{p{\linewidth}}}]{
                \texttt{\tiny{[P1$\rangle$GM]}}
                \texttt{To determine whether the argument is deductively valid or invalid, let's break down the logical structure of the premises and the conclusion.} \\
\\ 
\texttt{**1. Symbolize the Statements:**} \\
\texttt{Let:} \\
\texttt{*   \$C(x)\$: \$x\$ is a classmate of Williams.} \\
\texttt{*   \$W(x)\$: \$x\$ is a workmate of Leslie.} \\
\texttt{*   \$G(x)\$: \$x\$ is a great{-}grandfather of Joseph.} \\
\texttt{*   \$T(x)\$: \$x\$ is a cousin of Tony.} \\
\\ 
\texttt{**2. Analyze the Premises:**} \\
\\ 
\texttt{*   **Premise 1:** "Being a classmate of Williams is necessary for not being a workmate of Leslie."} \\
\texttt{    *   Logical form: If not \$W\$, then \$C\$.} \\
\texttt{    *   Symbolic: \$\textbackslash{}neg W(x) \textbackslash{}rightarrow C(x)\$.} \\
\texttt{    *   Contrapositive: \$\textbackslash{}neg C(x) \textbackslash{}rightarrow W(x)\$.} \\
\\ 
\texttt{*   **Premise 2:** "Being a classmate of Williams is necessary for not being a great{-}grandfather of Joseph."} \\
\texttt{    *   Logical form: If not \$G\$, then \$C\$.} \\
\texttt{    *   Symbolic: \$\textbackslash{}neg G(x) \textbackslash{}rightarrow C(x)\$.} \\
\texttt{    *   Contrapositive: \$\textbackslash{}neg C(x) \textbackslash{}rightarrow G(x)\$.} \\
\\ 
\texttt{*   **Premise 3:** "Every cousin of Tony is a workmate of Leslie and a great{-}grandfather of Joseph."} \\
\texttt{    *   Logical form: If \$T\$, then (\$W\$ and \$G\$).} \\
\texttt{    *   Symbolic: \$T(x) \textbackslash{}rightarrow (W(x) \textbackslash{}land G(x))\$.} \\
\\ 
\texttt{**3. Analyze the Conclusion:**} \\
\texttt{*   **Conclusion:** "Everyone who is a cousin of Tony is a classmate of Williams, too."} \\
\texttt{    *   Logical form: If \$T\$, then \$C\$.} \\
\texttt{    *   Symbolic: \$T(x) \textbackslash{}rightarrow C(x)\$.} \\
\\ 
\texttt{**4. Evaluate Validity:**} \\
\texttt{We need to see if the conclusion (\$T \textbackslash{}rightarrow C\$) necessarily follows from the premises.} \\
\\ 
\texttt{From Premise 3, we know that for any person \$x\$ who is a cousin of Tony (\$T\$), they are both a workmate of Leslie (\$W\$) and a great{-}grandfather of Joseph (\$G\$).} \\
\texttt{So, if \$T(x)\$ is true, then \$W(x)\$ is true and \$G(x)\$ is true.} \\
\\ 
\texttt{Now look at Premise 1 and Premise 2:} \\
\texttt{*   Premise 1 says: \$\textbackslash{}neg W \textbackslash{}rightarrow C\$. This tells us what happens if someone is *not* a workmate. It does **not** tell us what happens if someone *is* a workmate. (This is the fallacy of denying the antecedent if we tried to infer \$\textbackslash{}} \\
            }
        }
    }
    & & \\ \\

    \theutterance \stepcounter{utterance}  
    & & & \multicolumn{2}{p{0.3\linewidth}}{
        \cellcolor[rgb]{0.9,0.9,0.9}{
            \makecell[{{p{\linewidth}}}]{
                \texttt{\tiny{[GM$|$GM]}}
                \texttt{Target: invalid} \\
            }
        }
    }
    & & \\ \\

    \theutterance \stepcounter{utterance}  
    & & & \multicolumn{2}{p{0.3\linewidth}}{
        \cellcolor[rgb]{0.9,0.9,0.9}{
            \makecell[{{p{\linewidth}}}]{
                \texttt{\tiny{[GM$|$GM]}}
                \texttt{game\_result = LOSE} \\
            }
        }
    }
    & & \\ \\

    \theutterance \stepcounter{utterance}  
    & & \multicolumn{4}{p{0.6\linewidth}}{
        \cellcolor[rgb]{0.95,0.95,0.95}{
            \makecell[{{p{\linewidth}}}]{
                \texttt{\tiny{[P1$\leftrightarrow$P]}}
                \texttt{Post{-}episode Reflexion (not sent to GM):} \\
\texttt{New plan: I failed to complete the task because I stopped mid{-}sentence while explaining the logical fallacy (affirming the consequent) and never provided the final required answer option. My analysis was correct: the premises establish that \$\textbackslash{}neg W \textbackslash{}rightarrow C\$ and \$\textbackslash{}neg G \textbackslash{}rightarrow C\$, but the argument assumes \$W\$ and \$G\$ (from Premise 3) to conclude \$C\$, which is a logical error. Next time, I will ensure that after identifying the fallacy, I explicitly state the final verdict ("invalid") as the answer, rather than leaving the response incomplete. I will structure the response to first briefly outline the logical form, identify the specific error, and then immediately output the correct option.} \\
            }
        }
    }
    & \\ \\

\end{supertabular}
}

\end{document}
